\section{Introduction}
\label{sec:intro}

High-frequency bus service is vulnerable to bunching: a delayed vehicle faces
more waiting passengers and longer dwell, while its follower faces fewer,
causing an initially small headway perturbation to grow
~\citep{rezazada2024review}. Holding an early bus at selected stops can restore
spacing, but it also adds time for passengers already on board
~\citep{delgado2012holding}. Analytical headway control and schedule-reliability
control address different operational goals
~\citep{daganzo2009headway,xuan2011dynamic}; RL can represent richer stochastic
and multi-line interactions
~\citep{alesiani2018reinforcement,wang2020real,wang2023multiobjective}.

Learning through exploration on an operating fleet is generally unacceptable.
Offline RL can use historical transitions, but its policy is vulnerable outside
logged action support. Online learning in a simulator supplies adaptive
experience, but an inexpensive model can differ from the target in travel
times, dwell propagation, demand events, and the physical duration between
asynchronous decisions. H2O RL combines immutable
target replay with interaction in an imperfect simulator~\citep{niu2022h2o};
H2O+ generalizes this construction across offline and online base learners
~\citep{niu2025h2oplus}. The transit question is therefore not only whether
simulator interaction improves an offline policy, but whether the resulting
behavior remains service-complete when transferred back to the target
simulator.

Bus holding makes this distinction consequential. Stop decisions are
asynchronous, several buses can request control at the same physical time, and
exact zero hold is much more common than a positive duration. More importantly,
an objective that integrates waiting and in-vehicle passenger time does not by
itself ensure that journeys finish before a fixed evaluation horizon. During
development we observed precisely this objective--constraint mismatch: lower
generalized passenger time could coexist with a completion-rate loss large
enough to fail a preregistered service boundary. The final policy therefore
separates learning from deployment. A frozen H2O+ actor supplies the state
dependent action, while a deterministic completion-safety reserve modifies
only a bounded set of low-load tail-bus states with remaining demand.

The resulting evidence chain is summarized in \Cref{tab:evidence_chain}. The
target is a SUMO microsimulation of 12
directional services and 389 scheduled vehicle trips. The offline replay
contains 3,284,718 chronological transitions from 240 complete episodes. The
online environment is a causal event-driven simulator with matched timetable,
capacity, passenger-journey, state, action, and reward contracts. The final
candidate and direct parent share five frozen terminal-quarter actors and differ
only in the last deterministic completion-safety transform; lineage controls
retain their own frozen actors. Candidate design used prior development
outcomes, but the final candidate, direct parent, analysis, and ten-block formal
environment were fixed before formal outcomes were accessed.

\begin{table*}[t]
\centering
\small
\caption{Separation of learning, adaptive development, and formal confirmation.
Only the final column supplies confirmatory evidence.}
\label{tab:evidence_chain}
\begin{tabular}{>{\raggedright\arraybackslash}p{0.29\textwidth}
                >{\raggedright\arraybackslash}p{0.29\textwidth}
                >{\raggedright\arraybackslash}p{0.32\textwidth}}
\toprule
\textbf{Frozen learning evidence} & \textbf{Adaptive development} &
\textbf{Fresh formal confirmation} \\
\midrule
Target-SUMO replay; calibrated online simulator; five IQL-initialized H2O+
actors; 25,000 simulator events; fixed endpoint checkpoints. &
Terminal actor taper and deterministic hold transforms were developed over
successive preregistered rounds. Development outcomes informed the final
completion-safety reserve and therefore are not confirmatory. &
The fixed candidate and fixed direct parent were evaluated on 30 untouched
SUMO cells: ten independent blocks crossed with three demand levels. Exact
paired block tests used all five training seeds within each block. \\
\bottomrule
\end{tabular}
\end{table*}

Formal confirmation comprised 1,350 new SUMO episodes and 45 fixed physical
behaviors. Relative to its direct behavior parent, the reserve-augmented policy
reduced generalized passenger time by 92.12 s per departed passenger
($3.71\%$), increased completion by 0.00466, and reduced unfinished passengers
by 77.22 on average. The primary endpoint improved and the completion
noninferiority condition passed in all ten independent blocks. The result is
not attributed to density-ratio correction: held-out ratio diagnostics did not
generalize, and the sealed actor training used unit simulator weights.

This work makes four contributions. First, it defines a chronological
finite-horizon SMDP representation that serializes simultaneous
controls while charging exact generalized passenger time once per physical
interval. Second, it provides an audited multi-policy target dataset and a
causal online simulator under a common transit contract. Third, it combines a
stabilized H2O+ actor with an explicit, state-conditioned completion-safety
reserve whose direct effect can be isolated against an identical parent.
Fourth, it supplies a preregistered formal comparison on fresh environment
blocks, with the independent unit, noninferiority margin, and claim boundary
fixed before outcome access. The formal claim concerns the final reserve
relative to its identical direct parent in this simulator scenario family;
\Cref{sec:discussion:limitations} collects the limits on broader interpretation.
